\section{近期模型的 Scaling 实践}

\subsection{Llama 3}

\begin{figure}[H]
\centering
\includegraphics[width=0.95\textwidth]{slides-images/slide-025.jpg}
\caption{Llama 3 的 isoFLOP 分析：最优 token/parameter 比例约为 39:1\protect\footnotemark}
\end{figure}
\footnotetext{Slides 第26页。}

Llama 3 复现了 isoFLOP 分析，得到约 \textbf{39:1} 的最优 token/parameter 比例——高于 Chinchilla 的 20:1，但远低于 MiniCPM 的 192:1。

\begin{figure}[H]
\centering
\includegraphics[width=0.95\textwidth]{slides-images/slide-026.jpg}
\caption{Llama 3 的 Scaling-to-Benchmark 预测：通过拟合 sigmoid 函数，将 NLL 映射到下游任务准确率\protect\footnotemark}
\end{figure}
\footnotetext{Slides 第27页。}

Llama 3 的另一个有趣贡献是尝试将 scaling 从 log-likelihood 扩展到\textbf{下游任务准确率}。他们拟合了 sigmoid 函数，将 NLL（负对数似然）映射到 MMLU 等 benchmark 的分数，并成功预测了 Llama 3 405B 的下游性能。

\subsection{Hunyuan Large}

\begin{figure}[H]
\centering
\includegraphics[width=0.95\textwidth]{slides-images/slide-027.jpg}
\caption{Hunyuan Large 的 isoFLOP 分析：得到 96:1 的 data-to-active-parameter 比例\protect\footnotemark}
\end{figure}
\footnotetext{Slides 第28页。}

Hunyuan Large（混元大模型）同样复现了 Chinchilla 式分析，得到 \textbf{96:1} 的比例。由于架构差异（可能使用了 MoE 等），不同论文的比例自然会有所不同。

\subsection{MiniMax-01：Architecture Scaling}

\begin{figure}[H]
\centering
\includegraphics[width=0.95\textwidth]{slides-images/slide-028.jpg}
\caption{MiniMax-01：比较 Softmax attention、Lightning attention（线性）和 Hybrid attention 的 scaling 行为\protect\footnotemark}
\end{figure}
\footnotetext{Slides 第29页。}

MiniMax-01 展示了 scaling laws 的另一种应用：\textbf{架构选择验证}。他们比较了标准 Softmax attention、线性 Lightning attention 和混合模型的 scaling 曲线，发现三者在 compute-matched 下性能基本一致，从而为使用线性注意力实现长上下文提供了 scaling 层面的理论支撑。

\begin{knowledgebox}{用 Scaling Law 比较架构}
在 Mamba、DeltaNet 等线性注意力的研究论文中，经常看到类似的 compute-matched scaling 比较。MiniMax-01 的独特之处在于它不仅仅是学术论文中的小规模比较，而是在接近实际生产规模的模型上进行了验证。
\end{knowledgebox}

\subsection{案例研究总结}

\begin{figure}[H]
\centering
\includegraphics[width=0.95\textwidth]{slides-images/slide-029.jpg}
\caption{所有案例研究的 Scaling 策略对比\protect\footnotemark}
\end{figure}
\footnotetext{Slides 第30页。}

\begin{importantbox}{Scaling 实践中的共同模式}
\begin{itemize}
    \item \textbf{Chinchilla 复现}是最一致、最可靠的 scaling 分析——所有团队都做了，且结果非常干净
    \item \textbf{超参数 scaling}（学习率、batch size）总是更嘈杂，不同团队使用不同策略：MuP（Cerebras、MiniCPM）或直接拟合 scaling law（DeepSeek）
    \item \textbf{固定宽高比}后仅调整总模型规模，是处理架构 scaling 的通用做法
    \item \textbf{WSD 学习率调度}已被广泛采用，因其大幅降低了 data scaling 分析的成本
\end{itemize}
\end{importantbox}

\subsection{本章小结}

从 Chinchilla 原始论文的 20:1 到 Llama 3 的 39:1、Hunyuan 的 96:1，不同团队在不同条件下得到不同的最优比例。20:1 不是一个稳定的常数，但 isoFLOP 分析的\textbf{方法论}本身非常可靠且可复现。所有认真做 scaling 的团队都在某种程度上重复了 Chinchilla 的分析。

%% ========== Part 6: MuP 理论 ==========

